% ch9_d §9.8–9.9 (书页 324–328 / p338–p342)
%--C-TAIL--
%JOIN%
$\exp(-1/\omega_{H}\tau)$，也就是说，仿佛这个系统再也无法被冷却到某个固定温度 $T_{0}=\hbar/k\tau$ 之下。

上面的分析针对的是电子气的自由能和磁矩。系统的其他一些可观测性质，例如电导率和热导率，在强磁场中也会表现出类似的振荡效应。对这类行为作形式上的分析要复杂得多——但它本质上依赖于费米能级处的有效态密度随磁场的变化。振荡的电导率称为 de Haas--Shubnikov 效应。

\section{磁光吸收}

在半导体中，Landau 能级可以用光学方法直接探测。例如考虑一个有效质量为 $m^{*}$ 的电子的简单抛物线形导带。在磁量子化图式 (9.61) 中，各能级位于以变量 $k_{z}$ 描出的一系列抛物线上，彼此相隔一个磁量子 $\hbar\omega_{H}$（图 172）。

%FIG{172}: {(a) 抛物线能带的磁量子化图式。(b) 态密度。}

分配给每条抛物线各段的状态数由式 (9.67) 给出。对磁量子数 $n$ 的每个取值，我们都有一个简单的一维态密度问题。把能量 $\mathcal{E}$ 以下所有抛物线的贡献加起来，便得到总态密度函数
\begin{equation*}
\mathcal{N}_{H}(\mathcal{E})=\frac{1}{4\pi^{2}}\left(\frac{2m^{*}}{\hbar^{2}}\right)^{3/2}\hbar\omega_{H}\sum_{n}\left\{\mathcal{E}-\left(n+\frac{1}{2}\right)\hbar\omega_{H}\right\}^{-1/2}.
\tag{9.94}
\end{equation*}

%FIG{173}: {半导体中具有两条空穴能带时的磁光跃迁。}

当 $\omega_{H}\to 0$ 时，这个公式正确地积分后正确地回到标准的抛物线态密度函数 (4.33)。但磁场会在每个 Landau 能级处产生 $\mathcal{E}^{-1/2}$ 型的强 van Hove 奇异性（§2.5），这与 §9.7 中关于 de Haas--van Alphen 效应的一般论证相一致。

这类能级之间的跃迁将遵从选择定则 $\Delta n=\pm 1$；这不过是回旋共振（§9.2）的另一种描述方式。但态密度中的磁结构也足够宽，足以在其他能带的红外跃迁谱中被观察到，正如 §8.5 中所讨论的。这就是基本的磁光效应（magneto-optical effect）。

要解释实际的观测结果，还必须计及价带态的量子化，其 $\omega_{H}$ 和 $m^{*}$ 取不同的值（图 173）。对于一个在没有磁场时本来会发生的典型强跃迁，我们可以取矩阵元 (8.40) 与 $H$ 无关，并施加磁选择定则 $\Delta n=0$。于是联合态密度 (8.73) 便取与式 (9.94) 相同的形式；当 $\omega$ 与带间共振频率
\begin{equation*}
\hbar\omega_{n}=\mathcal{E}_{G}+\left(n+\frac{1}{2}\right)\hbar\left(\omega_{e}+\omega_{h}\right)
\tag{9.95}
\end{equation*}
相重合时，就会出现 $(\omega-\omega_{n})^{-1/2}$ 型的奇异性。不过在实际中，这些奇异性会因碰撞而展宽。

这样，这一现象就为测绘电子与空穴的回旋频率 $\omega_{e}$ 和 $\omega_{h}$ 提供了非常精确的手段，其能量范围覆盖带隙两侧。在一般的各向异性情形下，若同时存在价带简并以及能量对 $|\mathbf{k}|$ 的非抛物线型依赖，磁光谱可能非常复杂，但它提供了关于晶体电子能带结构的极为有价值的证据。

\section{磁击穿}

在本章中，我们讨论的一直是强磁场中的电子。特别地，我们研究了电子在被散射之前已回旋许多圈时出现的种种现象。我们已经看到，这时采用一种新的量子化图式是合宜的：在这一图式中，磁能级把 Bloch 图式中分立的各个允许态都吸收了进来。当然，如果存在散射，磁能级本身也会展宽，但只要 $\omega_{H}\tau\gg 1$，这一效应就可以忽略。

但是，我们提出的这种对每条轨道都满足式 (9.70) 的量子化图式，并不是在真正极其巨大的磁场中应当使用的终极图式。§9.6 中所用的准经典理论依赖于 §6.4 的等效哈密顿量原理；\emph{只有当我们能够忽略带间跃迁时，它才是有效的。}如果磁场非常大，这类跃迁就必然会发生——如果你愿意这么说，是通过隧穿发生的，正如 Zener 效应（§6.8）中那样——于是 \emph{Onsager} 图式 (9.70) 就会失效。

要理解这一效应，最容易的办法是从磁场极高的一端出发：在那里，电子波函数本质上就是带电粒子在磁场中自由运动的波函数，而晶格的晶体势仅作为微扰。把我们限制在垂直于磁场的二维空间内，我们在 $(k_{x},k_{y})$ 平面上便有一条圆形轨道。

现在引入如下形式的微扰
\begin{equation*}
\mathcal{V}(x)=\sum_{g}\mathcal{V}_{g}e^{igx}
\tag{9.96}
\end{equation*}
——即一系列相距 $2\pi/g$ 的平面。当轨道穿过区边界时，也就是说，当它沿 $x$ 方向的有效波矢 $k_{x}$ 等于 $\pm\frac{1}{2}g$ 时，就有可能发生 Bragg 衍射。轨道不再（比如说）沿 $AB$ 继续前进，而可能转到 $AC$ 方向，如此等等。

%FIG{174}: {(a) 磁场中的自由电子轨道。(b) 在晶格的周期势中，轨道在区边界处被重新连接；但在强磁场下，轨道可能跳回自由电子的路径。}

如果我们增大微扰的强度，那么 $A$ 处的轨道将在能量上分裂开，而 $AC$ 这条路径将成为优先的。此时电子便在通常的重复区图式中沿一条寻常的（开）轨道运动。图 174(a) 中圆的 $B$ 部分这时已被连接成费米面的一个单独分支，被完全分开地走完。

如果我们这时增大磁场，我们将回到圆形轨道图式。电子不再总是沿 $AC$ 前进，而是可能击穿能隙，或者说“隧穿”通过倒格子空间中分隔两条轨道的区域，发觉自己绕着 $B$ 运动了。

事实上，我们可以用 Zener 击穿的公式 (6.65) 来估计这个跃迁概率。电场 $E$ 能够引起穿过能隙 $\mathcal{E}_{\mathrm{gap}}$ 的隧穿，如果
\begin{equation*}
\frac{eEa\mathcal{E}^{0}}{(\mathcal{E}_{\mathrm{gap}})^{2}}>1,
\tag{9.97}
\end{equation*}
其中 $\mathcal{E}^{0}$ 是电子在能隙处的动能——实际上就是费米能 $\mathcal{E}_{F}$——而 $a$ 是相应的晶格间距。

现在，在重复区图式中趋向 $A$ 的电子具有速度
\begin{equation*}
v\sim\frac{\hbar k_{F}}{m},
\tag{9.98}
\end{equation*}
其中 $k_{F}$ 是费米半径。这在区边界处并不严格成立，因为等能面会因能隙的存在而发生畸变，但我们正假设这个畸变很小。

电子以这一速度横越磁场运动，会产生一个 Lorentz 力，它等同于一个强度为
\begin{equation*}
E\sim\frac{vH}{c}
\tag{9.99}
\end{equation*}
、方向与 $v$ 成直角的电场。如果式 (9.97) 得到满足，也就是说，如果参数
\begin{equation*}
\frac{ev\mathbf{H}}{c}\,\frac{a\mathcal{E}^{0}}{(\mathcal{E}_{\mathrm{gap}})^{2}}\approx\frac{e\hbar}{mc}\,H\,\frac{k_{F}a\mathcal{E}_{F}}{(\mathcal{E}_{\mathrm{gap}})^{2}}
\tag{9.100}
\end{equation*}
大于 1，这个场就能引起隧穿。由 (3.3)，乘积 $k_{F}a$ 是一个量级为 1 的数。利用 (9.4)，磁击穿（magnetic breakdown）的条件便成为
\begin{equation*}
\frac{\hbar\omega_{H}\mathcal{E}_{F}}{(\mathcal{E}_{\mathrm{gap}})^{2}}>1.
\tag{9.101}
\end{equation*}

这个条件远不如（我们不妨这么说）磁轨道的间隔必须大于能隙这样的条件那样苛刻。对某些金属，在 100 kG 量级的磁场中它就变得重要起来，并且会在 de Haas--van Alphen 效应以及其他现象中引起各种不同面积的新轨道的出现。
